% ============================================================
%  RELATÓRIO TÉCNICO — Template SBC
%  Como usar (Overleaf, sem instalar nada):
%   1. Acesse overleaf.com > New Project > Templates > procure "SBC"
%      (template "SBC Conferences" — ele já vem com o sbc-template.sty).
%   2. Apague o conteúdo do main.tex de exemplo e cole ESTE arquivo.
%   3. Suba suas imagens (prints do jogo) em Upload e ajuste os \includegraphics.
%   4. Preencha tudo marcado com  %% PREENCHER %%  e clique em Recompile.
%
%  Procure por  PREENCHER  para achar rapidinho o que falta.
% ============================================================

\documentclass[12pt]{article}

\usepackage{sbc-template}
\usepackage{graphicx,url}
\usepackage[utf8]{inputenc}
\usepackage[brazil]{babel}

\sloppy

\title{Flappy Capivara: uma reinterpretação 3D do Flappy Bird em OpenGL}

\author{Eduardo Melo de Carvalho\inst{1}, Maria Eduarda Farias Gomes\inst{1}}

\address{Departamento de Computação -- Centro de Ciências da Natureza\\
  Universidade Federal do Piauí (UFPI) -- Teresina, PI -- Brasil\\
  Matrículas: %% PREENCHER: matricula do Eduardo, matricula da Maria Eduarda
  \email{eduardo.melo@ufpi.edu.br, maria.farias@ufpi.edu.br}
}

\begin{document}

\maketitle

\begin{abstract}
This paper presents \emph{Flappy Capivara}, a 3D reinterpretation of the
classic 2D game \emph{Flappy Bird} developed in C++ with OpenGL and GLUT.
The player controls a capybara that must fly through gaps between pipes,
avoiding the pipes and the ground, while an allied bee, driven by an
autopilot artificial intelligence, flies ahead and navigates the gaps on
its own. The project applies core Computer Graphics techniques: a 3D
perspective camera, lighting and shading, textures, visibility through
depth buffer and back-face culling, 3D model loading, collision detection
with bounding volumes, and a reactive artificial intelligence. We describe
the implemented methods, present results, and discuss the main challenges.
\end{abstract}

\begin{resumo}
Este artigo apresenta o \emph{Flappy Capivara}, uma reinterpretação em 3D
do jogo clássico 2D \emph{Flappy Bird}, desenvolvido em C++ com OpenGL e
GLUT. O jogador controla uma capivara que deve voar pelas brechas entre
canos, desviando dos canos e do chão, enquanto uma abelha aliada, guiada por
uma IA de piloto automático, voa à frente e atravessa as brechas sozinha.
O projeto aplica técnicas fundamentais de
Computação Gráfica: câmera 3D em perspectiva, iluminação e sombreamento,
texturas, visibilidade por \emph{buffer} de profundidade e \emph{back-face
culling}, carregamento de modelos 3D, detecção de colisão por volumes
envolventes e uma inteligência artificial reativa. Descrevemos os métodos
implementados, apresentamos os resultados e discutimos os principais
desafios.
\end{resumo}

% ============================================================
\section{Introdução}
% ============================================================

Jogos clássicos em duas dimensões são um excelente ponto de partida para o
estudo de Computação Gráfica, pois permitem explorar técnicas de
renderização tridimensional sobre uma jogabilidade simples e bem conhecida.
Seguindo essa ideia, este trabalho reinterpreta o jogo \emph{Flappy Bird}
em três dimensões, substituindo o pássaro original por uma \textbf{capivara}
--- animal símbolo da fauna brasileira --- e dando ao cenário um visual 3D.

No \emph{Flappy Capivara}, a câmera observa a cena em perspectiva, enquanto
a jogabilidade permanece no plano 2D (eixos $X$ e $Y$), atendendo ao
requisito de um jogo 3D com jogabilidade ao menos em 2D. A capivara cai
continuamente por ação da gravidade e sobe a cada batida de asa (acionada
pelo jogador). O objetivo é atravessar o maior número possível de pares de
canos sem colidir.

Além de reproduzir a mecânica original, o jogo adiciona elementos
inexistentes no clássico: uma \textbf{abelha aliada com inteligência
artificial} (piloto automático),
um \textbf{cenário com vegetação e profundidade} (\emph{parallax}),
\textbf{partículas}, \textbf{áudio} e \textbf{modos de dificuldade}. Os
objetivos do trabalho
são: (i) construir um jogo 3D funcional em OpenGL/GLUT; (ii) aplicar
detecção de colisões, inteligência artificial e os recursos de iluminação,
sombreamento, textura e visibilidade; e (iii) integrar bibliotecas
complementares de modelos 3D, áudio e fontes.

% ============================================================
\section{Materiais}
% ============================================================

O jogo foi desenvolvido em \textbf{C++} e compilado com o \emph{clang}
(macOS). As principais ferramentas e bibliotecas, com suas finalidades,
são:

\begin{itemize}
  \item \textbf{OpenGL} --- API de renderização gráfica (geometria,
        iluminação, textura, \emph{buffer} de profundidade).
  \item \textbf{GLUT} --- criação de janela, contexto de renderização e
        tratamento de eventos de teclado, \emph{mouse} e \emph{idle}.
  \item \textbf{Assimp} --- carregamento de modelos 3D nos formatos
        \texttt{.obj} e \texttt{.glb}.
  \item \textbf{stb\_image} --- leitura de imagens PNG usadas como textura.
  \item \textbf{stb\_truetype} + fonte \textbf{Pixelify Sans} --- renderização
        do texto do HUD (placar e mensagens).
  \item \textbf{miniaudio} --- saída de áudio (os efeitos sonoros são
        sintetizados em tempo real pelo próprio código).
  \item \textbf{Blender} --- modelagem e exportação de \emph{assets} 3D
        (por exemplo, o cano com borda separada, em formato \texttt{.glb}).
\end{itemize}

Os \emph{assets} 3D foram obtidos no repositório livre \emph{poly.pizza}
(modelos em formato \texttt{.obj}, de licença livre): a \textbf{capivara},
a \textbf{asa de morcego}, a \textbf{abelha} aliada, o \textbf{tufo de
grama} e os dois modelos de \textbf{árvore} do cenário de fundo. O
\textbf{cano} foi remodelado no \textbf{Blender}, com corpo e borda como
peças separadas, e exportado em \texttt{.glb}. Quanto às texturas, a
capivara usa uma imagem PNG (\emph{base color}), enquanto a textura do chão
é \textbf{gerada proceduralmente} no próprio código.

% ============================================================
\section{Métodos}
% ============================================================

\subsection{Janela, câmera e projeção}
A janela e o contexto são criados com GLUT em \emph{double buffering}. A
projeção é em \textbf{perspectiva} (\texttt{gluPerspective}), o que dá a
sensação de profundidade. A câmera é posicionada com \texttt{gluLookAt},
levemente afastada no eixo $Z$, observando o plano de jogo.

\subsection{Iluminação e sombreamento}
A cena é iluminada por uma \textbf{fonte de luz direcional}
(\texttt{GL\_LIGHT0} com o quarto componente da posição igual a zero),
simulando o \textbf{Sol} com raios paralelos. A iluminação combina as três
componentes do modelo de reflexão: \textbf{difusa} (tom levemente quente),
\textbf{ambiente} (preenchimento frio, para as sombras não ficarem pretas) e
\textbf{especular}, com o material recebendo um brilho controlado por um
expoente de \emph{shininess}. O sombreamento é suave (\emph{Gouraud}), com
normais por vértice geradas no carregamento dos modelos. O recurso
\texttt{GL\_COLOR\_MATERIAL} permite definir o material dos objetos
diretamente pela cor, e \texttt{GL\_NORMALIZE} mantém a iluminação correta
mesmo quando os modelos são escalados. As cores de alguns objetos são ainda
levemente \emph{dessaturadas} para harmonizar a paleta do jogo.

\subsection{Texturas}
A capivara é texturizada a partir de uma imagem PNG (\emph{base color})
carregada com \emph{stb\_image} e mapeada pelas coordenadas de textura
(\emph{UV}) do modelo. O chão usa uma \textbf{textura procedural} de grama,
\emph{gerada por código} (variações de verde e tons terrosos), aplicada com
repetição (\emph{tiling}, \texttt{GL\_REPEAT}) e filtragem \texttt{GL\_NEAREST}
para um aspecto pixelado. Modelos em \texttt{.glb}/\texttt{.obj} sem imagem de
textura são coloridos pelas cores difusas definidas em seus materiais.

\subsection{Visibilidade}
A visibilidade correta entre objetos é garantida pelo \textbf{\emph{buffer}
de profundidade} (\emph{z-buffer}, \texttt{GL\_DEPTH\_TEST}). Como recurso
adicional, habilita-se o \textbf{\emph{back-face culling}}
(\texttt{GL\_CULL\_FACE}), que descarta as faces traseiras dos objetos,
reduzindo o custo de desenho.

\subsection{Carregamento de modelos 3D}
A função de carregamento usa o Assimp para importar a malha, triangular as
faces e gerar normais suaves. Em seguida, calcula-se a \emph{caixa
envolvente} (\emph{bounding box}) de cada modelo para normalizar seu tamanho
e centro, padronizando a escala dos objetos no mundo.

\subsection{Canos: corpo e borda}
Cada par de canos é desenhado a partir de um modelo \texttt{.glb} criado no
Blender. Para encaixar a brecha em alturas variadas, o \textbf{corpo} do
cano (um cilindro liso) é esticado verticalmente, enquanto a \textbf{borda}
(\emph{rim}) é desenhada com escala \textbf{uniforme} na boca do cano. Essa
separação evita a distorção que ocorria quando o modelo inteiro era
esticado.

\subsection{Cenário: profundidade, vegetação e céu}
Para reforçar a percepção tridimensional, o cenário é composto em camadas.
Sobre o \textbf{chão} texturizado rolam \textbf{tufos de grama} 3D
distribuídos em várias \emph{fileiras} de profundidade, formando um gramado.
Ao \textbf{fundo}, \textbf{árvores} de dois modelos, com escala e
profundidade sorteadas, deslocam-se em \textbf{\emph{parallax}} --- mais
lentamente que os canos --- o que acentua a sensação de distância. Grama,
árvores e canos seguem a estratégia de \textbf{reciclagem}
(\emph{object pooling}): os elementos que saem pela esquerda reaparecem à
direita com novos parâmetros, evitando alocação de memória durante a
partida. O \textbf{céu} é desenhado como um \emph{quad} em projeção
ortográfica com \textbf{degradê} vertical, no lugar de uma cor chapada. Por
fim, a geometria de cada modelo é compilada uma única vez em
\textbf{\emph{display lists}}, evitando reenviar os vértices à GPU a cada
quadro --- otimização importante dado o grande número de instâncias de grama
e árvores desenhadas por quadro.

\subsection{Física e controle}
O movimento vertical usa integração de Euler semi-implícita: a cada quadro,
a velocidade recebe a aceleração da gravidade e a posição é atualizada pela
velocidade, usando o tempo decorrido ($\Delta t$) para independência da taxa
de quadros. A batida de asa (tecla \emph{espaço} ou clique) aplica um
impulso vertical e dispara a animação das asas.

\subsection{Detecção de colisões}
A colisão usa \textbf{caixas envolventes alinhadas aos eixos} (AABB): a
capivara é aproximada por uma caixa e cada cano por outra; há colisão quando
as caixas se sobrepõem simultaneamente nos eixos $X$ e $Y$. O chão é tratado
como um limite inferior em $Y$. Qualquer colisão encerra a partida. O teste é
isolado em uma função pura (\texttt{sobreposicaoAABB}), validada por testes
automatizados. O projeto também implementa o cálculo de distância entre
centros, $D=\sqrt{(x_1-x_0)^2+(y_1-y_0)^2}$, base da colisão por esferas
envolventes, igualmente coberto por testes.

\subsection{Inteligência artificial}
Uma \textbf{abelha aliada} é controlada por uma IA de \textbf{piloto
automático} que \emph{joga sozinha}. Ela permanece numa posição fixa à
frente da capivara e, a cada quadro: (i) identifica o próximo cano em seu
caminho; (ii) adota como alvo o \emph{centro da brecha} desse cano; e (iii)
aplica a mesma física da capivara --- gravidade somada a uma ``batida de
asa'' automática sempre que está abaixo do alvo. Assim, a aliada navega
pelas brechas como um jogador e funciona como um \textbf{guia}, mostrando o
caminho seguro à frente. Trata-se de uma \textbf{agente reativa}: percebe o
ambiente (a posição da próxima brecha) e age por regras simples, o que a
torna fácil de explicar e robusta em todas as dificuldades.

\subsection{Modos de dificuldade}
O jogo oferece \textbf{três modos}, escolhidos pelas teclas 1, 2 e 3 na tela
inicial. No \textbf{Fácil}, a brecha entre os canos e a velocidade são fixas.
No \textbf{Médio}, a brecha \emph{encolhe} conforme a pontuação --- até um
\emph{limite mínimo} que mantém o jogo jogável --- e a velocidade é fixa. No
\textbf{Difícil}, a brecha encolhe \emph{e} a velocidade aumenta com a
pontuação (também com tetos). Assim, a mesma base de jogo cobre desde um
treino tranquilo até um desafio progressivo.

\subsection{Demais recursos}
O jogo inclui ainda: um \textbf{sistema de partículas} (poeira ao passar o
cano, brilho ao pontuar e explosão ao colidir) desenhado com
\emph{billboards} e \emph{blending}; \emph{feedback} de impacto na colisão
com \textbf{tremor de tela} (\emph{screen shake}) e \textbf{flash}; um
\textbf{HUD} em projeção ortográfica, com o título em fonte
\emph{Pixelify Sans} rasterizada por \emph{stb\_truetype} (atlas de textura,
com gradiente vertical e contorno) além do placar e das mensagens; e uma
\textbf{máquina de estados do jogo} (início, jogando, fim de jogo).
Quanto ao \textbf{áudio}, os efeitos sonoros (pulo, ponto e colisão) são
\textbf{sintetizados em tempo real} --- ondas quadradas e ruído com envelope
de decaimento, no estilo \emph{8-bit} --- por um \emph{callback} do
\emph{miniaudio}, sem depender de arquivos de som, o que mantém o projeto
autocontido e livre de questões de direitos autorais.

% ============================================================
\section{Resultados e Discussão}
% ============================================================

O jogo executa de forma fluida, com a capivara respondendo ao comando do
jogador, os canos rolando e reciclando, o placar incrementando a cada cano
ultrapassado e a abelha aliada navegando as brechas à frente da capivara.
As Figuras~\ref{fig:inicio}, \ref{fig:jogo} e \ref{fig:gameover} mostram,
respectivamente, a tela inicial com a escolha de dificuldade, o jogo em
execução e a tela de fim de jogo.

\begin{figure}[ht]
  \centering
  \includegraphics[width=.7\textwidth]{tela_inicial.png}
  \caption{Tela inicial: título e escolha de dificuldade (teclas 1/2/3).}
  \label{fig:inicio}
\end{figure}

\begin{figure}[ht]
  \centering
  \includegraphics[width=.7\textwidth]{jogando.png}
  \caption{Jogo em execução: capivara, canos com borda, abelha aliada e placar.}
  \label{fig:jogo}
\end{figure}

\begin{figure}[ht]
  \centering
  \includegraphics[width=.7\textwidth]{game_over.png}
  \caption{Tela de fim de jogo com a pontuação final.}
  \label{fig:gameover}
\end{figure}

A combinação de iluminação direcional, sombreamento suave e textura confere
volume aos objetos, enquanto o \emph{z-buffer} e o \emph{back-face culling}
garantem a visibilidade correta. As camadas de grama e árvores em
\emph{parallax}, somadas ao céu em degradê, dão profundidade ao cenário; as
partículas, o tremor de tela e o áudio sintetizado reforçam o
\emph{feedback} de jogo. O uso de \emph{display lists} manteve a execução
fluida mesmo com muitas instâncias de vegetação. A colisão por AABB
mostrou-se simples e eficaz para a jogabilidade. Por fim, a IA de piloto
automático da abelha aliada demonstra o requisito de inteligência artificial
de forma marcante, pois ela efetivamente \emph{joga} o jogo, navegando as
brechas à frente do jogador.

% ============================================================
\section{Conclusão}
% ============================================================

O desenvolvimento do \emph{Flappy Capivara} permitiu aplicar, de forma
integrada, os principais conceitos de Computação Gráfica estudados na
disciplina. Entre os \textbf{desafios} enfrentados destacam-se: a animação
correta das asas, resolvida com a separação das duas metades; a distorção da
borda do cano, resolvida modelando corpo e borda como peças separadas no
Blender e escalando-as de forma independente; e o ajuste das colisões e da
inteligência artificial para uma jogabilidade equilibrada.

Como \textbf{trabalhos futuros}, sugerem-se: animações esqueletais
importadas, novos tipos de obstáculos e de aliados/inimigos, fases com
cenários distintos e \emph{ranking} de pontuação.

% ============================================================
\section*{Referências}
% ============================================================
\begin{thebibliography}{9}

\bibitem{opengl} OpenGL. \emph{OpenGL Documentation}.
  Disponível em: \url{https://www.opengl.org/}.

\bibitem{glut} Kilgard, M. J. \emph{The OpenGL Utility Toolkit (GLUT)}.

\bibitem{assimp} Open Asset Import Library (Assimp).
  \url{https://www.assimp.org/}.

\bibitem{stb} Barrett, S. \emph{stb single-file public domain libraries}.
  \url{https://github.com/nothings/stb}.

\bibitem{miniaudio} \emph{miniaudio}. \url{https://miniaud.io/}.

\bibitem{blender} Blender Foundation. \emph{Blender}.
  \url{https://www.blender.org/}.

\bibitem{polypizza} Poly Pizza. \emph{Free 3D models}.
  \url{https://poly.pizza/}.

\bibitem{pixelify} Google Fonts. \emph{Pixelify Sans}.
  \url{https://fonts.google.com/specimen/Pixelify+Sans}.

\bibitem{sbc} Sociedade Brasileira de Computação (SBC).
  \emph{Templates para artigos}. \url{https://www.sbc.org.br/}.

\end{thebibliography}

\end{document}
